\documentclass[11pt]{article}
\usepackage[T1]{fontenc}
\usepackage[utf8]{inputenc}
\usepackage{lmodern}
\usepackage{microtype}
\usepackage[margin=1in]{geometry}
\usepackage{amsmath,amssymb,amsthm,mathtools}
\usepackage{booktabs,longtable,array}
\usepackage{enumitem}
\usepackage{xcolor}
\usepackage[hidelinks]{hyperref}
\usepackage[nameinlink,noabbrev]{cleveref}

\newtheorem{theorem}{Theorem}[section]
\newtheorem{lemma}[theorem]{Lemma}
\newtheorem{proposition}[theorem]{Proposition}
\newtheorem{corollary}[theorem]{Corollary}
\newtheorem{definition}[theorem]{Definition}
\newtheorem{remark}[theorem]{Remark}
\newtheorem{construction}[theorem]{Construction}

\newcommand{\N}{\mathbb N}
\newcommand{\Z}{\mathbb Z}
\newcommand{\Q}{\mathbb Q}
\newcommand{\R}{\mathbb R}
\newcommand{\ipra}{\mathsf{iPRA}}
\newcommand{\PR}{\mathsf{PR}}
\newcommand{\Gam}{\boldsymbol\Gamma}
\newcommand{\Eq}{\mathrel{=_{\!\mathsf R}}}
\newcommand{\Le}{\mathrel{\le_{\!\mathsf R}}}
\newcommand{\Pos}{\mathsf{Pos}}
\newcommand{\RLC}{\mathsf{RLC}}
\newcommand{\BM}{\mathsf{BM}}
\newcommand{\code}[1]{\texttt{#1}}
\newcommand{\ceil}[1]{\left\lceil #1\right\rceil}
\newcommand{\abs}[1]{\left|#1\right|}
\newcommand{\prodj}[2]{\prod_{j=#1}^{#2}}

\title{A Finitary Constructive Bohr--Mollerup Theorem in Intuitionistic Primitive Recursive Arithmetic}
\author{Research report prepared for Vladimir Reshetnikov}
\date{15 November 2026}

\begin{document}
\maketitle

\begin{abstract}
The classical Bohr--Mollerup theorem characterizes the gamma function as the
unique positive log-convex function satisfying $f(1)=1$ and
$f(x+1)=x f(x)$.  Neither arbitrary functions on the real line nor an
unqualified real-completeness principle are objects of first-order
intuitionistic primitive recursive arithmetic.  We therefore give a precise
finitary reformulation.  Reals are primitive-recursively generated nested
rational intervals with an explicit rate, functions on positive rationals are
fixed primitive-recursive interval presentations with local moduli, and
log-convexity is replaced by an algebraic rational-weight condition
\[
 F\!\left(\frac{(s-j)x+jy}{s}\right)^s
 \le F(x)^{s-j}F(y)^j.
\]
A gamma presentation is constructed directly from Euler's finite products.
On the fundamental strip $0<q\le1$ the two brackets
\[
 A_n(q)=\frac{n!\,n^q}{q(q+1)\cdots(q+n)},\qquad
 B_n(q)=\frac{(n-1)!\,n^q}{q(q+1)\cdots(q+n-1)}
\]
are nested and satisfy $0\le B_n(q)-A_n(q)\le1/n$.  Rational powers are
implemented by bounded bisection of integer roots, so these brackets yield an
explicit primitive-recursive real name.  The recurrence extends the name to
all positive arguments.  A finite two-sided squeeze proves uniqueness, with
an explicit value of $n$ for every requested precision.  Every witness and
modulus is primitive recursive; all induction can be reduced to
quantifier-free induction.  No logarithm, integral, unbounded search,
countable choice, excluded middle for arbitrary propositions, or general
Cauchy-completeness axiom is used.  The report ends with a module-by-module
Lean formalization plan tied to the existing \code{ConciseIPRA} development.
\end{abstract}

\tableofcontents

\section{What is proved}

The exact classical statement quantifies over arbitrary functions
$F:(0,\infty)\to(0,\infty)$.  Such a higher-type quantifier is not part of the
first-order language of $\ipra$.  A mathematically honest formulation must
therefore say how a real and a function are presented, and it must distinguish
an internal arithmetic theorem from the external notation used to read a
regular Cauchy presentation as a real number.

This report proves two closely related results.

\begin{enumerate}[label=\textbf{(\Alph*)},leftmargin=2.5em]
\item \emph{Internal rational-core theorem schema.}  There is a fixed
primitive-recursive interval presentation $\Gam_{\Q}$ on positive rational
codes.  For every fixed primitive-recursive presentation $F$ satisfying the
constructive Bohr--Mollerup hypotheses, $\ipra$ proves, for every positive
rational code $q$ and precision $k$, that $F(q)$ and $\Gam_{\Q}(q)$ agree to
precision $2^{-k}$.  ``For every fixed presentation'' is a theorem schema,
which is the appropriate first-order reading: the defining primitive-recursive
functions of $F$ are function symbols of the particular definitional
extension.

\item \emph{Completion corollary.}  If the presentation of $F$ also supplies
a primitive-recursive local uniform-continuity modulus, then both $F$ and
$\Gam_{\Q}$ extend canonically from positive rationals to regular positive
reals, and the extensions are equal.  This is again witnessed precision by
precision, not by an appeal to an impredicative set of all real numbers.
\end{enumerate}

The algebraic convexity hypothesis is equivalent to ordinary log-convexity in
classical analysis once positivity and continuity are available, but it is
substantially better suited to $\ipra$: its individual instances contain only
finite products, integer powers, rational arithmetic, and order.

\section{Arithmetic coding and regular real names}
\label{sec:coding}

\subsection{Rationals}

A nonnegative rational is coded by a pair $(a,b)$ with $a\in\N$ and
$0<b\in\N$; signed rationals add a sign bit.  Equality and order are decided
by cross multiplication.  Addition, multiplication, reciprocal on a supplied
positive denominator, finite products, factorial, powers, reduction by a
bounded gcd computation, and maximum or minimum of a finite list are all
primitive recursive.  No appeal to a field object is needed inside the
arithmetic theory.

We write rational notation in the paper.  Every displayed rational identity
can be expanded into an integer identity, and every rational inequality into
an integer inequality after multiplication by explicitly positive
denominators.

\subsection{Nested interval names}

\begin{definition}[regular interval name]
A \emph{regular interval name} is a pair of primitive-recursive rational
functions $L_X,U_X:\N\to\Q$ such that, for every $k$,
\begin{align*}
 L_X(k)&\le L_X(k+1)\le U_X(k+1)\le U_X(k),\\
 U_X(k)-L_X(k)&\le2^{-k}.
\end{align*}
Its center is $c_X(k)=(L_X(k)+U_X(k))/2$.
\end{definition}

A convenient arithmetic equality predicate is
\begin{equation}
 X\Eq Y\quad:\Longleftrightarrow\quad
 \forall k\;\abs{c_X(k+2)-c_Y(k+2)}\le2^{-k}.
 \label{eq:real-equality}
\end{equation}
Shifting the precision by two makes reflexivity, symmetry, and transitivity
immediate from the triangle inequality and regularity.  Other standard regular
Cauchy conventions are definitionally interchangeable with this one.  The
order relation is treated similarly:
\begin{equation}
 X\Le Y\quad:\Longleftrightarrow\quad
 \forall k\; c_X(k+2)\le c_Y(k+2)+2^{-k}.
 \label{eq:real-order}
\end{equation}
All predicates at a fixed precision are decidable rational inequalities.

\begin{definition}[strong positivity]
A name $X$ is \emph{strongly positive} if its data include integers $p,s$ with
$2^{-s}\le L_X(p)$.  By nesting, this is then a permanent positive lower
bound.  We write $\Pos(X)$.
\end{definition}

This is the constructive meaning of ``positive'' used below.  It is stronger
than merely denying $X\le0$, and it supplies exactly the witness required to
form reciprocals and cancel a positive factor.

The usual arithmetic operations on regular names are obtained by interval
arithmetic followed by a precision shift.  For example, if $X,Y$ have known
bounds on the precision range in use, the product interval is the minimum and
maximum of the four endpoint products.  A running-intersection normalization
\[
 \widehat L(k)=\max_{i\le k}L(i),\qquad
 \widehat U(k)=\min_{i\le k}U(i)
\]
turns any compatible sequence of shrinking enclosures into a nested one.  All
of these transformations are primitive recursive.

\subsection{Presented functions and the first-order status}

For the internal theorem, a candidate function on positive rationals is a
fixed tuple of primitive-recursive functions
\[
 L_F(q,k),\quad U_F(q,k),\quad P_F(q),\quad M_F(\alpha,\beta,k),
\]
where $q,\alpha,\beta$ are rational codes.  The first two give a regular name,
$P_F$ gives a strong-positivity witness, and $M_F$ is a local modulus:
whenever $\alpha\le x,y\le\beta$ and
$\abs{x-y}\le2^{-M_F(\alpha,\beta,k)}$, the values agree within $2^{-k}$.
The theorem on rational arguments does not need $M_F$; it is needed only for
the completion corollary.

Because $F$ is fixed, its primitive-recursive algorithms occur as ordinary
function symbols in a definitional extension of $\ipra$.  The theorem has the
form
\[
 \BM(F)\longrightarrow
 \forall q\,\forall k\,
 \bigl(q>0\longrightarrow \mathsf{EqAt}_k(F(q),\Gam_{\Q}(q))\bigr).
 \label{eq:schema-shape}
\]
There is no hidden quantification over all functions.  Uniformity is supplied
metatheoretically by a primitive-recursive proof constructor taking the fixed
syntax of the presentation and its hypothesis proofs to the syntax of the
conclusion proof.

\section{The constructive convexity condition}

\begin{definition}[rational multiplicative convexity]
A strongly positive presented function $F$ satisfies $\RLC$ if for all
positive rational $x,y$, all $s\ge1$, and all $0\le j\le s$, with
\[
 z=\frac{(s-j)x+jy}{s},
\]
one has
\begin{equation}
 F(z)^s\Le F(x)^{s-j}F(y)^j.
 \label{eq:rlc}
\end{equation}
\end{definition}

The condition contains no logarithm and no irrational scalar.  Classically,
for positive continuous functions, it is equivalent to convexity of
$\log F$: one direction follows by exponentiating Jensen's inequality; the
other follows first at rational weights from \eqref{eq:rlc} and then at real
weights by density and continuity.

\begin{definition}[constructive Bohr--Mollerup data]
The assertion $\BM(F)$ consists of regularity and strong positivity of every
$F(q)$, together with
\begin{align}
 F(1)&\Eq1, \label{eq:normalization}\\
 F(x+1)&\Eq xF(x) \qquad(x>0), \label{eq:recurrence}\\
 \RLC(F)&. \label{eq:convexity-assumption}
\end{align}
For the real-domain corollary, the local modulus $M_F$ is also included.
\end{definition}

\section{Finite weighted arithmetic--geometric mean}
\label{sec:amgm}

The only convexity inequality needed for constructing gamma is a finite
weighted AM--GM inequality.

\begin{lemma}[finite AM--GM, constructive form]
For positive rationals $u_1,\dots,u_r$,
\[
 \left(\frac{u_1+\cdots+u_r}{r}\right)^r
 \ge u_1\cdots u_r.
\]
There is a primitive-recursive algebraic certificate of this inequality from
the input rational codes.
\end{lemma}

\begin{proof}
The two-variable case is the identity
\[
 \left(\frac{u-v}{2}\right)^2\ge0.
\]
If the result is known for $r$ entries, splitting $2r$ entries into two blocks,
applying the $r$-entry result to each block, and then the two-entry result to
their two means proves the $2r$-entry result.  Conversely, the $r$-entry result
implies the $(r-1)$-entry result: append the mean
$A=(u_1+\cdots+u_{r-1})/(r-1)$ as the last entry.  AM--GM for $r$ entries gives
$A^r\ge A u_1\cdots u_{r-1}$; the supplied positive lower bound for $A$ permits
cancellation.  To handle arbitrary $r$, choose the least power of two at least
$r$, prove that case by repeated doubling, and descend one step at a time.

Every loop is bounded by $r$ or by the binary length of $r$.  The certificate
consists only of additions, multiplications, squares, substitutions, and
cancellation by explicitly positive rationals.  Its correctness is proved by
quantifier-free induction on the finite construction tree.
\end{proof}

\begin{corollary}[two-point rational weight]
For $u,v>0$, $0\le a\le b$, and $b>0$,
\begin{equation}
 \left(\frac{(b-a)u+av}{b}\right)^b\ge u^{b-a}v^a.
 \label{eq:weighted-amgm}
\end{equation}
\end{corollary}

\begin{proof}
Apply finite AM--GM to a list containing $b-a$ copies of $u$ and $a$ copies
of $v$.
\end{proof}

When $t>0$ and $r=a/b\in[0,1]$, \eqref{eq:weighted-amgm} says, in the notation
of positive $b$th roots,
\begin{equation}
 t^r\le (1-r)+rt.
 \label{eq:concave-power}
\end{equation}
The notation $t^{a/b}$ is constructive: for rational $t$, it denotes the
unique positive $b$th root of $t^a$, represented by bounded bisection.  The
comparison in \eqref{eq:concave-power} can equivalently be checked after
raising both sides to the $b$th power, so no analytic root theorem is required
in the proof checker.

\section{Primitive-recursive rational powers}
\label{sec:roots}

For the Euler products below, only $n^{a/b}$ with $n\in\N_{>0}$ and
$0<a\le b$ is needed.  Put $T=n^a$.  The positive $b$th root lies in $[1,n]$.
At each bisection step, for a rational midpoint $m$, the decidable integer
comparison
\[
 \operatorname{num}(m)^b
 \le T\,\operatorname{den}(m)^b
\]
selects the lower or upper half.  After $h$ steps, the width is at most
$n2^{-h}$.  Thus, for any positive rational $\delta$, choosing the first
$h$ with $2^h\ge n/\delta$ gives bounds
\[
 \rho^-\le n^{a/b}\le\rho^+,
 \qquad \rho^+-\rho^-\le\delta.
\]
The search for $h$ is bounded by the bit lengths of the numerator and
denominator of $n/\delta$, hence is primitive recursive.  This construction
also proves monotonicity and uniqueness of positive integer roots by finite
integer-power comparisons.

\section{Euler brackets on the fundamental strip}
\label{sec:strip}

Fix a rational $q=a/b$ in lowest terms with $0<q\le1$.  For $n\ge1$ define
\begin{align}
 A_n(q)&=\frac{n!\,n^q}{q(q+1)\cdots(q+n)},
 \label{eq:An}\\
 B_n(q)&=\frac{(n-1)!\,n^q}{q(q+1)\cdots(q+n-1)}.
 \label{eq:Bn}
\end{align}
The relation between them is
\begin{equation}
 B_n(q)=\left(1+\frac qn\right)A_n(q).
 \label{eq:BoverA}
\end{equation}

\begin{lemma}[nested Euler brackets]
For $0<q\le1$ and $n\ge1$,
\begin{equation}
 A_n(q)\le A_{n+1}(q)\le B_{n+1}(q)\le B_n(q),
 \label{eq:nesting}
\end{equation}
and
\begin{equation}
 0\le B_n(q)-A_n(q)\le\frac1n.
 \label{eq:width}
\end{equation}
\end{lemma}

\begin{proof}
Direct cancellation gives
\begin{align*}
 \frac{A_{n+1}(q)}{A_n(q)}
 &=\frac{n+1}{n+1+q}\left(1+\frac1n\right)^q,\\
 \frac{B_{n+1}(q)}{B_n(q)}
 &=\frac{n}{n+q}\left(1+\frac1n\right)^q.
\end{align*}
Apply \eqref{eq:concave-power} first to $t=n/(n+1)$:
\[
 \left(\frac n{n+1}\right)^q\le1-\frac q{n+1}.
\]
Taking positive reciprocals and using the elementary identity
$(1-s)^{-1}\ge1+s$ for $0\le s<1$ gives
\[
 \left(1+\frac1n\right)^q
 \ge1+\frac q{n+1}=\frac{n+1+q}{n+1},
\]
so $A_{n+1}\ge A_n$.  Applying \eqref{eq:concave-power} to
$t=1+1/n$ gives
\[
 \left(1+\frac1n\right)^q\le1+\frac qn=\frac{n+q}{n},
\]
so $B_{n+1}\le B_n$.  The middle inequality in \eqref{eq:nesting} is
immediate from \eqref{eq:BoverA}.

Now $B_n\le B_1=1/q$, and hence $A_n\le1/q$.  From
\eqref{eq:BoverA},
\[
 B_n-A_n=\frac qn A_n\le\frac1n.
\]
Every step is an inequality between positive rationals after taking the
$b$th power; no limiting argument has occurred.
\end{proof}

\subsection{An explicit rational interval name}

Although $A_n(q)$ and $B_n(q)$ contain the algebraic number $n^q$, rational
endpoints are obtained by the bisection construction of \cref{sec:roots}.
Write
\[
 A_n=c_A n^q,\qquad B_n=c_B n^q,
\]
where $c_A,c_B$ are positive rationals.  Since $n^q\ge1$ and
$A_n,B_n\le1/q$, one has $c_A,c_B\le1/q$.

At requested precision $k$, take
\[
 n_k=2^{k+2},\qquad
 \delta_k=q\,2^{-k-3}.
\]
Compute rational root bounds $\rho_k^-\le n_k^q\le\rho_k^+$ with width at
most $\delta_k$, and put
\[
 \ell_k=c_A\rho_k^-,\qquad u_k=c_B\rho_k^+.
\]
Then
\begin{align*}
 u_k-\ell_k
 &\le (B_{n_k}-A_{n_k})+(c_A+c_B)\delta_k\\
 &\le2^{-k-2}+2\delta_k/q
 =2^{-k-1}.
 \label{eq:raw-width}
\end{align*}
For $i\le j$, nesting gives
$\ell_i\le A_{n_i}\le A_{n_j}\le B_{n_j}\le u_j$; by the same chain with
the indices reversed at the outer endpoints, every $\ell_i$ is below every
$u_j$.  Therefore
\[
 L_0(k)=\max_{i\le k}\ell_i,
 \qquad U_0(k)=\min_{i\le k}u_i
 \label{eq:strip-name}
\]
is a nested nonempty rational interval of width at most $2^{-k-1}$.  This is
our primitive-recursive name $\Gam_0(q)$ on $0<q\le1$.

By construction,
\begin{equation}
 A_n(q)\Le\Gam_0(q)\Le B_n(q)
 \qquad(n\ge1).
 \label{eq:gamma-strip-squeeze}
\end{equation}
Here an algebraic endpoint is itself read through its bisection name; all
precision instances reduce to rational inequalities.

\section{Extension to all positive rationals}
\label{sec:extension}

Every positive rational has a canonical decomposition
\[
 x=m+r,
 \qquad m=\ceil{x}-1\in\N,
 \qquad 0<r\le1.
\]
Let
\[
 P_m(r)=\prod_{j=0}^{m-1}(r+j),\qquad P_0(r)=1,
\]
and define
\begin{equation}
 \Gam_{\Q}(x)=P_m(r)\Gam_0(r).
 \label{eq:global-gamma}
\end{equation}
This is a primitive-recursive interval presentation: $m,r$, the finite
product, and the required precision shift for multiplication are all explicit.

\begin{proposition}[normalization, recurrence, and positivity]
For positive rational $x$,
\begin{align*}
 \Gam_{\Q}(1)&\Eq1,\\
 \Gam_{\Q}(x+1)&\Eq x\Gam_{\Q}(x),\\
 \Pos(\Gam_{\Q}(x))&.
\end{align*}
\end{proposition}

\begin{proof}
For $q=1$, the strip brackets are
$A_n(1)=n/(n+1)$ and $B_n(1)=1$, so
\eqref{eq:gamma-strip-squeeze} and the explicit width force
$\Gam_0(1)=1$.  The recurrence is immediate from
$P_{m+1}(r)=(m+r)P_m(r)$; the boundary case $r=1$ is the same factorial
identity.  Strong positivity follows from the explicit positive lower
endpoint of any one strip bracket and the positive finite product $P_m(r)$.
\end{proof}

\section{Convexity of the constructed gamma presentation}
\label{sec:gamma-convexity}

For $x>0$ rational and $n\ge1$, define the global Euler approximant
\begin{equation}
 E_n(x)=\frac{n!\,n^x}{x(x+1)\cdots(x+n)}.
 \label{eq:En}
\end{equation}

\begin{lemma}[explicit convergence from below]
Write $x=m+r$ with $0<r\le1$, and set
\[
 C_m=m!\left(1+m+\frac{m(m+1)}2\right).
\]
Then
\begin{equation}
 0\Le \Gam_{\Q}(x)-E_n(x)\Le\frac{C_m}{n}.
 \label{eq:En-tail}
\end{equation}
\end{lemma}

\begin{proof}
A direct rearrangement gives
\begin{equation}
 E_n(m+r)=P_m(r)A_n(r)
 \prod_{s=1}^{m}\frac{n}{n+r+s}.
 \label{eq:En-factor}
\end{equation}
The empty product is $1$.  Hence $E_n\le P_mA_n\le\Gam_{\Q}$.
Moreover,
\begin{align*}
 \Gam_{\Q}-E_n
 &=P_m(\Gam_0-A_n)
   +P_mA_n\left(1-\prod_{s=1}^{m}\frac{n}{n+r+s}\right).
\end{align*}
The first term is at most $P_m/n\le m!/n$.  For the second, use
$A_n\le1/r$, $P_m/r\le m!$, and the finite product inequality
$1-\prod_s(1-d_s)\le\sum_s d_s$:
\begin{align*}
 1-\prod_{s=1}^{m}\frac{n}{n+r+s}
 &\le\sum_{s=1}^{m}\frac{r+s}{n+r+s}\\
 &\le\frac{mr+m(m+1)/2}{n}\\
 &\le\frac{m+m(m+1)/2}{n}.
\end{align*}
Combining the two estimates gives \eqref{eq:En-tail}.  This is a finite
calculation with an explicit error term.
\end{proof}

\begin{lemma}[each Euler approximant is rationally log-convex]
If $z=((s-j)x+jy)/s$, then
\begin{equation}
 E_n(z)^s\le E_n(x)^{s-j}E_n(y)^j.
 \label{eq:En-rlc}
\end{equation}
\end{lemma}

\begin{proof}
The numerator factors cancel because $sz=(s-j)x+jy$.  For every
$h=0,\dots,n$, weighted AM--GM gives
\[
 (z+h)^s\ge(x+h)^{s-j}(y+h)^j.
\]
Multiplying these finitely many inequalities makes the denominator on the
left at least the corresponding denominator on the right, which reverses in
the quotient and yields \eqref{eq:En-rlc}.
\end{proof}

\begin{proposition}[constructive rational multiplicative convexity]
The presentation $\Gam_{\Q}$ satisfies $\RLC$.
\end{proposition}

\begin{proof}
Fix $x,y,s,j$ and $z=((s-j)x+jy)/s$.  By \eqref{eq:En-tail}, each of
$E_n(x),E_n(y),E_n(z)$ approaches the corresponding gamma name with a stated
$C/n$ error.  Choose an explicit rational upper bound $H\ge1$ for all six
quantities; for example, use the strip bound $\Gam_0(r)\le1/r$, the finite
shift products, and then enlarge to cover the finitely many $E_n$.
For $0\le u,v\le H$,
\begin{equation}
 \abs{u^d-v^d}\le dH^{d-1}\abs{u-v}
 \label{eq:power-lipschitz}
\end{equation}
follows from factorization of $u^d-v^d$.  A second telescoping factorization
gives an explicit analogous bound for $u^{s-j}v^j$.

Given precision $2^{-k}$, take $n$ larger than the maximum of the three tail
constants times the resulting integer-power Lipschitz constant times
$2^{k+2}$.  Then \eqref{eq:En-rlc} implies
\[
 \Gam_{\Q}(z)^s
 \le \Gam_{\Q}(x)^{s-j}\Gam_{\Q}(y)^j+2^{-k}.
\]
Since $k$ was arbitrary, this is exactly the regular-name order
\eqref{eq:real-order}.  The choice of $n$ is an explicit primitive-recursive
expression in the input codes and $k$.
\end{proof}

\section{The finite uniqueness squeeze}
\label{sec:uniqueness}

Let $F$ be any presented positive function satisfying $\BM(F)$.  Fix
$q=a/b$ in lowest terms with $0<q\le1$, and write
\[
 P_n(q)=q(q+1)\cdots(q+n-1).
\]
Repeated use of normalization and recurrence gives
\begin{equation}
 F(n)=(n-1)!,\qquad F(n+1)=n!,\qquad
 F(n+q)=P_n(q)F(q).
 \label{eq:recurrence-iterated}
\end{equation}

\begin{lemma}[upper squeeze]
For every $n\ge1$,
\begin{equation}
 F(q)^b\Le B_n(q)^b.
 \label{eq:upper-squeeze}
\end{equation}
\end{lemma}

\begin{proof}
Because
\[
 n+q=\frac{(b-a)n+a(n+1)}b,
\]
rational multiplicative convexity and \eqref{eq:recurrence-iterated} give
\begin{align*}
 P_n(q)^bF(q)^b
 &=F(n+q)^b\\
 &\le F(n)^{b-a}F(n+1)^a\\
 &=((n-1)!)^b n^a.
\end{align*}
Division by the explicitly positive $P_n(q)^b$ gives exactly
\eqref{eq:upper-squeeze}, because
$B_n(q)^b=((n-1)!)^b n^a/P_n(q)^b$.
\end{proof}

\begin{lemma}[lower squeeze]
For every $n\ge1$,
\begin{equation}
 A_n(q)^b\Le F(q)^b.
 \label{eq:lower-squeeze}
\end{equation}
\end{lemma}

\begin{proof}
The identity
\[
 n+1=\frac{a(n+q)+(b-a)(n+q+1)}b
\]
and $\RLC$ yield
\begin{align*}
 (n!)^b=F(n+1)^b
 &\le F(n+q)^aF(n+q+1)^{b-a}\\
 &=F(n+q)^b(n+q)^{b-a}\\
 &=P_n(q)^bF(q)^b(n+q)^{b-a}.
\end{align*}
Thus
\[
 F(q)^b\ge\frac{(n!)^b}{P_n(q)^b(n+q)^{b-a}}.
\]
Since $n^a\le(n+q)^a$, the right side is at least
\[
 \frac{(n!)^b n^a}{P_n(q)^b(n+q)^b}=A_n(q)^b.
\]
\end{proof}

The constructed gamma value satisfies the same squeeze by
\eqref{eq:gamma-strip-squeeze}.  The interval for the $b$th powers shrinks at
an explicit rate:
\begin{lemma}[power-gap bound]
\begin{equation}
 0\le B_n(q)^b-A_n(q)^b
 \le \frac{b}{nq^{b-1}}.
 \label{eq:power-gap}
\end{equation}
\end{lemma}

\begin{proof}
Factor the difference:
\[
 B_n^b-A_n^b=(B_n-A_n)
 \sum_{i=0}^{b-1}B_n^{b-1-i}A_n^i.
\]
By \eqref{eq:width}, $B_n-A_n\le1/n$, and by the proof of that lemma,
$0\le A_n\le B_n\le1/q$.  There are $b$ summands, each at most
$q^{1-b}$.
\end{proof}

\begin{theorem}[constructive Bohr--Mollerup theorem on positive rationals]
Let $F$ be a fixed primitive-recursive regular interval presentation on
positive rationals satisfying $\BM(F)$.  Then
\[
 \forall q\in\Q_{>0}\qquad F(q)\Eq\Gam_{\Q}(q).
\]
For every $q$ and precision $k$, a witnessing derivation uses only finitely
many recurrence instances and two instances of $\RLC$, at a value of $n$
computed primitively from $q$ and $k$.
\end{theorem}

\begin{proof}
First let $0<q=a/b\le1$.  The upper and lower squeeze lemmas show that both
$F(q)^b$ and $\Gam_0(q)^b$ lie between $A_n(q)^b$ and $B_n(q)^b$ for every
$n$.  Given precision $2^{-k}$, choose, for example,
\[
 n\ge 2^{k+2}\,b\,q^{1-b}.
 \label{eq:explicit-n-unique}
\]
Then \eqref{eq:power-gap} proves that the two $b$th powers agree to the
required precision.  Hence
$F(q)^b\Eq\Gam_0(q)^b$.

Positive integer powers are injective on strongly positive regular names.
Indeed,
\[
 F(q)^b-\Gam_0(q)^b
 =(F(q)-\Gam_0(q))
 \sum_{i=0}^{b-1}F(q)^{b-1-i}\Gam_0(q)^i.
\]
The second factor has an explicit positive lower bound from the supplied
positivity witnesses, so it has a constructively defined reciprocal and can
be cancelled.  Thus $F(q)\Eq\Gam_0(q)$.

For arbitrary $x=m+r>0$ with $0<r\le1$, iterate the recurrence finitely:
\[
 F(x)=P_m(r)F(r)=P_m(r)\Gam_0(r)=\Gam_{\Q}(x).
\]
All choices made in this proof are displayed bounded computations.  In
particular, no proof by contradiction or nonconstructive extraction of a
limit is used.
\end{proof}

\section{A primitive-recursive local modulus and the real-domain corollary}
\label{sec:completion}

It remains to justify that the rational presentation extends canonically to
regular positive reals.  We give an explicit local modulus rather than
invoking a general continuity theorem.

Fix positive rationals $0<\alpha\le\beta$, let
$M=\ceil{\beta}$, and put $C=C_M$.  For rational
$x,y\in[\alpha,\beta]$, choose $n$ so that $2C/n\le\varepsilon/3$.
The denominator
\[
 D_n(x)=\prod_{j=0}^{n}(x+j)
\]
satisfies $D_n(x)\ge\alpha n!$, and a telescoping product estimate gives
\begin{equation}
 \abs{D_n(x)-D_n(y)}
 \le(n+1)(\beta+n)^n\abs{x-y}.
 \label{eq:den-lipschitz}
\end{equation}
For an integer $T\ge1$, Bernoulli's finite inequality gives
\[
 \left(1+\frac{n-1}{T}\right)^T\ge n,
\]
so $n^{1/T}\le1+(n-1)/T$.  Consequently, if
$\abs{x-y}\le1/T$, then, using $n^t\le n^{\ceil\beta}$ on the interval,
\begin{equation}
 \abs{n^x-n^y}
 \le n^{\ceil\beta}\frac{n-1}{T}.
 \label{eq:power-modulus}
\end{equation}
Combining \eqref{eq:den-lipschitz}, \eqref{eq:power-modulus}, and the lower
bound for $D_n$ yields
\begin{equation}
 \abs{E_n(x)-E_n(y)}\le\frac{K_{\alpha,\beta,n}}{T},
 \label{eq:En-modulus}
\end{equation}
where the rational constant
\begin{equation}
 K_{\alpha,\beta,n}
 =n^{\ceil\beta}
 \left(
   \frac{n-1}{\alpha}
   +\frac{(n+1)(\beta+n)^n}{\alpha^2 n!}
 \right)
 \label{eq:Kconstant}
\end{equation}
is primitive recursively computable.  Choose $T$ with
$K_{\alpha,\beta,n}/T\le\varepsilon/3$.  By the two tail estimates,
\[
 \abs{\Gam_{\Q}(x)-\Gam_{\Q}(y)}\le\varepsilon.
\]
Setting $\varepsilon=2^{-k}$ and replacing $1/T$ by the next smaller dyadic
produces a primitive-recursive local modulus
$M_{\Gam}(\alpha,\beta,k)$.

\begin{corollary}[regular-real Bohr--Mollerup theorem]
The rational presentation $\Gam_{\Q}$ has a canonical extension $\Gam$ to
strongly positive regular real names.  If a presented function $F$ has a
primitive-recursive local modulus and satisfies normalization, recurrence, and
rational multiplicative convexity, then
\[
 \forall X>0\qquad F(X)\Eq\Gam(X).
\]
Moreover $\Gam$ itself satisfies the three conditions.
\end{corollary}

\begin{proof}
Given a regular name $X$ and a compact positive rational interval known to
contain it, take the canonical rational center $x_t=c_X(t)$ at a precision
selected by $M_{\Gam}$.  The values $\Gam_{\Q}(x_t)$ form a regular name by
the modulus; this defines $\Gam(X)$ without choice.  Independence of the
chosen approximation follows from the same modulus.  Normalization,
recurrence, and $\RLC$ pass from rational approximants to real names using
finite product moduli.

For uniqueness, select $t$ large enough for both the modulus of $F$ and that
of $\Gam$.  Then
\[
 F(X)\approx F(x_t)=\Gam_{\Q}(x_t)\approx\Gam(X)
\]
with each outer error below one third of the requested precision.  Since the
rational-core equality is already proved and every precision is explicit, the
regular names are equal.
\end{proof}

\section{Why the proof belongs to \texorpdfstring{$\ipra$}{iPRA}}
\label{sec:ipra-audit}

The preceding argument is constructive in a stronger sense than merely
avoiding an obvious appeal to excluded middle.  Its computational inventory
is as follows.

\begin{longtable}{>{\raggedright\arraybackslash}p{0.30\textwidth}
                  >{\raggedright\arraybackslash}p{0.63\textwidth}}
\toprule
Component & Primitive-recursive realization and induction used \\
\midrule
\endhead
Rational arithmetic & Integer pairs, cross multiplication, bounded gcd,
finite sums and products; all equations and inequalities are decidable.\\
Regular real name & Two PR endpoint functions; regularity and equality are
universal closures of quantifier-free endpoint inequalities.\\
Finite AM--GM & A bounded proof tree built by doubling and descent; correctness
by quantifier-free induction on its code.\\
Rational powers & A prescribed number of dyadic bisection steps; branch choice
is a decidable integer-power comparison.\\
Euler brackets & Factorial, powers, rising products, and rational arithmetic;
monotonicity is checked after a finite integer power.\\
Gamma name & Running maximum and minimum of finitely many rational endpoints;
no completeness axiom.\\
Convexity & A finite AM--GM product for each $E_n$, followed by an explicitly
chosen tail index.\\
Uniqueness & Two convexity instances, finitely many recurrence instances, and
the explicit bound \eqref{eq:explicit-n-unique}.\\
Real extension & Canonical rational centers and the explicit local modulus
\eqref{eq:Kconstant}; no choice sequence.\\
\bottomrule
\end{longtable}

The induction axioms of the usual first-order presentation of $\ipra$ apply to
quantifier-free formulas.  This suffices because every invariant above is a
primitive-recursive equality or inequality on codes.  Universal statements
are obtained by generalizing a free numerical parameter after the
quantifier-free induction has been performed.  Finite sequences are decoded
by primitive-recursive projections, so their correctness predicates are also
decidable.

No step uses any of the following:
\begin{itemize}[leftmargin=2em]
\item a logarithm or exponential function on arbitrary reals;
\item an integral representation of gamma;
\item a least-upper-bound or general Cauchy-completeness axiom;
\item countable or dependent choice;
\item an unbounded search for a convergence stage;
\item excluded middle for an arbitrary proposition;
\item a nonconstructive claim that every positive real has a chosen root.
\end{itemize}
Only decidable case distinctions on integer comparisons occur.

\section{Lean formalization roadmap}
\label{sec:lean}

The repository's existing \code{ConciseIPRA} development already supplies the
right object theory.  In particular, its usual presentation has a type
\code{PR k} of primitive-recursive function symbols, terms \code{PTm}, axioms
\code{PRAAx}, and the theorem \code{PRA.qfInd} for quantifier-free induction.
The file \code{BiInterpretation.lean} supplies
\code{pra\_theorem\_iff}, translating closed theorems between the usual
presentation and the concise scalar theory.  The cleanest formalization has
two layers.

\subsection{Layer I: executable arithmetic and semantic correctness}

The following modules are proposed.

\begin{description}[style=nextline,leftmargin=2.7cm]
\item[\code{BM/Code.lean}]
Pairing, finite sequences, signed rational codes, cross-multiplied equality and
order, dyadics, factorial, rising products, bounded maxima/minima.

\item[\code{BM/IntervalReal.lean}]
Regular interval names, equality/order at a precision, strong positivity,
interval arithmetic, running-intersection normalization, and precision-shift
lemmas.

\item[\code{BM/AMGM.lean}]
The proof-tree datatype for finite AM--GM, its evaluator, and correctness of
the two-point weighted instance.

\item[\code{BM/Root.lean}]
Bounded bisection for positive integer roots, invariants, width bound, and
rational-power comparisons after clearing the root by an integer power.

\item[\code{BM/EulerBracket.lean}]
Definitions of $A_n,B_n$, exact ratio identities, monotonicity, the $1/n$
width theorem, and rational endpoint generation.

\item[\code{BM/GammaQ.lean}]
The strip name, running intersection, shift decomposition $x=m+r$, global
gamma name, normalization, recurrence, and positivity.

\item[\code{BM/EulerConvex.lean}]
$E_n$, the factorization \eqref{eq:En-factor}, tail bound, finite weighted
convexity, and passage to gamma with explicit power-error bounds.

\item[\code{BM/Squeeze.lean}]
The two finite convexity applications, power-gap bound, positive-power
cancellation, and rational-core uniqueness.

\item[\code{BM/Completion.lean}]
The explicit modulus \eqref{eq:Kconstant}, extension from rational centers, and
real-domain uniqueness.
\end{description}

These modules may use Lean's ordinary natural-number semantics to debug the
algorithms and prove their external correctness.  The supplied Python program
\code{code/bm\_bounds.py} is a small independent executable model of the same
finite arithmetic.

\subsection{Layer II: object-level \texorpdfstring{$\ipra$}{iPRA} proofs}

\begin{description}[style=nextline,leftmargin=2.7cm]
\item[\code{BM/PRDefs.lean}]
Compile every algorithm from Layer I into a value of the existing
\code{PR k} syntax.  Prove its defining equations using
\code{PRA.precZero}, \code{PRA.precSucc}, composition, and projection.

\item[\code{BM/QFGraphs.lean}]
Define quantifier-free graph formulas for rational operations, bisection,
finite products, and interval predicates.  Prove functionality and totality
for the fixed PR terms.

\item[\code{BM/ObjectLemmas.lean}]
Mirror the finite algebra lemmas as derivations
$\Gamma \vdash[\mathsf{PRAAx}] \varphi$.  Use \code{PRA.qfInd}; avoid importing any
classical real-analysis theorem.

\item[\code{BM/ObjectTheorem.lean}]
Assemble the closed formula expressing \eqref{eq:schema-shape} for a fixed
candidate presentation and build the derivation of every precision instance.
At this layer, the candidate hypotheses are formulas in the context, not
meta-level assumptions silently discharged.

\item[\code{BM/Translate.lean}]
Apply \code{pra\_theorem\_iff} to transport the final closed theorem to the
concise native scalar theory.  Add \code{\#print axioms} checks and a no-
\code{sorryAx} audit script.
\end{description}

\subsection{Suggested milestone order}

A low-risk route to a complete formalization is:
\begin{enumerate}
\item certify rational coding and $A_n,B_n$ numerically;
\item formalize the nesting and width theorem;
\item obtain the strip and global gamma names;
\item prove the two uniqueness squeezes, initially as an abstract ordered
semiring theorem;
\item instantiate that theorem for interval names;
\item prove convexity of the constructed gamma using $E_n$;
\item add the local modulus and completion layer;
\item compile the PR algorithms and replay the proof inside \code{PRAAx};
\item translate through the existing bi-interpretation and audit axioms.
\end{enumerate}
The archive contains a commented Lean scaffold with the intended declarations
and dependency graph.  It is an outline, not a completed machine-checked
formalization.

\section{Constructive-content audit of the classical proof}

The familiar textbook proof contains a genuinely finitary core.  The present
reformulation isolates it:

\begin{center}
\begin{tabular}{ll}
\toprule
Classical presentation & Constructive replacement \\
\midrule
$\log F$ is convex & integer-power inequality \eqref{eq:rlc} \\
real interpolation parameter & rational pair $(j,s)$ \\
Euler limit exists & explicit nested rational brackets \\
general completeness of $\R$ & a named PR interval algorithm \\
``let $n\to\infty$'' & choose the displayed $n(q,k)$ \\
continuity inferred abstractly & supplied or derived PR local modulus \\
positive real & positive name with a rational lower witness \\
function quantifier & schema for each fixed PR presentation \\
\bottomrule
\end{tabular}
\end{center}

This is not merely the classical proof with intuitionistic vocabulary.  The
change from logarithmic convexity to rational multiplicative convexity, and
the replacement of an unnamed limit by the $A_n$--$B_n$ interval algorithm,
are the two essential proof-theoretic modifications.

\section{Further research questions}

\begin{enumerate}
\item Can the local-modulus hypothesis on a candidate be derived inside the
chosen real representation from recurrence and $\RLC$, with a uniform
primitive-recursive bound, rather than supplied as part of the presentation?

\item What is the smallest fragment of the \code{ConciseIPRA} native scalar
calculus needed for the finite squeeze?  The uniqueness proof itself appears
to need substantially less machinery than the construction and completion
parts.

\item Can the finite AM--GM certificate generator be replaced by a shorter
binomial or sum-of-squares certificate that materially reduces the final
object-level proof term?

\item Is there a direct gamma presentation whose rational endpoint sequence is
already nested, eliminating the running maximum/minimum normalization and
thereby simplifying the Lean object proof?

\item Can one characterize a wider class of factorial-like functions in
$\ipra$ by replacing $x$ in the recurrence with a fixed positive
primitive-recursive coefficient function and adapting the same finite squeeze?

\item How small can a standalone proof checker and certificate for a fixed
precision instance of the theorem be made under the code-size metrics used in
the \code{ConciseIPRA} project?
\end{enumerate}

\appendix
\section{Exact certificate checked by the companion program}

The companion program accepts a positive rational $x$ and a requested number
of bits $k$.  It performs the canonical decomposition $x=m+r$, chooses an
internal strip precision, computes $n=2^{K+2}$, bisects the integer root
$n^{\operatorname{num}(r)}$ to bound $n^r$, and emits rational endpoints.
The verifier recomputes:
\begin{itemize}
\item the root-power inequalities;
\item the exact $c_A,c_B$ coefficients;
\item the shifted interval endpoints;
\item the requested width bound;
\item the two integer-power inequalities certifying $A_n\le A_{n+1}$ and
$B_{n+1}\le B_n$ at the selected $n$.
\end{itemize}
Verification uses Python's exact \code{Fraction} arithmetic only.  Decimal
numbers printed by the command-line interface are display-only.

\section{Dependency map to the existing repository}

The proposed object-level development imports only the syntactic and proof
infrastructure already present in the \code{ConciseIPRA} tree:
\[
 \code{ConciseIPRA.PR}
 \longrightarrow \code{BM.ObjectLemmas}
 \longrightarrow \code{BM.ObjectTheorem},
\]
followed by
\[
 \code{ConciseIPRA.BiInterpretation}
 \longrightarrow \code{BM.Translate}.
\]
The existing usual presentation provides a symbol for every primitive
recursive definition, its defining equations, successor axioms, and
quantifier-free induction.  The existing bi-interpretation proves that closed
theorems of this presentation are exactly the translations of closed theorems
of the concise scalar theory.  Thus the mathematical work in this report can
be formalized first in the more readable usual language without losing the
connection to the constrained native calculus.

\begin{thebibliography}{9}
\bibitem{DLMF}
NIST Digital Library of Mathematical Functions,
\emph{Chapter 5: Gamma Function}, especially the Bohr--Mollerup
characterization in Section 5.5(iv).

\bibitem{BishopBridges}
Errett Bishop and Douglas Bridges,
\emph{Constructive Analysis}, Springer, 1985.

\bibitem{MathlibGamma}
The mathlib community,
\emph{Mathlib.Analysis.SpecialFunctions.Gamma.BohrMollerup}, Lean source and
API documentation.  The present proof does not depend on its classical real
analysis; it is useful as a theorem-name and regression reference.

\bibitem{ConciseIPRA}
Vladimir Reshetnikov,
\emph{Concise presentations of intuitionistic primitive recursive
arithmetic}, the \code{Logic/PrimitiveRecursiveArithmetic/ConciseIPRA}
subtree of the ProveIt repository, including \code{PR.lean} and
\code{BiInterpretation.lean}.
\end{thebibliography}

\end{document}
